单模无源Kerr环腔在驱动旋转系中满足
\(\dot a=[i(\Delta-g|a|^2)-\kappa/2]a+\sqrt{\kappa_{\rm ex}}s_{\rm in}\)，其中 \(g>0\)，\(\kappa=\kappa_{\rm ex}+\kappa_0>0\)，\(\kappa_0\ge0\)，\(n=|a|^2\)，入射光子流 \(P_{\rm in}=|s_{\rm in}|^2\)，输出振幅 \(s_{\rm out}=s_{\rm in}-\sqrt{\kappa_{\rm ex}}a\)。
(1) 求稳态光子数的无量纲三次方程，说明有效失谐的Kerr符号约定。
(2) 求产生双稳态的失谐阈值、折叠位置和尖点条件，解释多根数目与稳定态数目不等价。
(3) 在复振幅空间推导稳态的线性化Jacobian和稳定性判据，区分三分支的动力学性质。
(4) 从时变场方程独立证明光子流输入输出平衡，求强驱动与零非线性极限，并辨析实际滞回中噪声的作用。
